\clearpage
\section{Podsumowanie analiza kosztów}

    \paragraph{}
        Zgodnie z rodziałem 2.7 Analiza kosztów ksiąrzki \cite{galinski} mojego promotora. 
        Wykonałem wykres zależności odwrotności obciążenia ciągu od obciążenia powierzchni.
        napoczątku przed wykonaniem obliczen sprawdziłem ile czy wykonałem poprawnie obliczenia.
        Szacujac koszty dla zakonczonej juz produkcji 737-300 z lat powstania raportu na podstawie ktorego.
        Postała methoda wyszło mi rok 1970: 20980423 USD oraz w roku 1986: 16368673 USD co wydaje się dos prawdopodobne.
        Ponieważ smaloty były sprzedawane w tamtym okresie za 40 - 46,5 mln USD. Przyjmuje średnią cenne 43,5 mln USD i daje 2.66 spredu.

        \begin{itemize}
            \item Rok 1970:
                \begin{itemize}
                    \item Wsparcie prac badawczo rozwojowych: 381526 USD
                    \item Cena silnika i awioniki: 9041976 USD
                    \item Materiały: 16095760 USD
                    \item Próby w locie: 23137136 USD
                    \item Koszt inzynierów po przliczeniu stawek za roboczo godziny: 914704 USD
                    \item Koszt narzędzi po przliczeniu stawek za roboczo godziny: 31959501 USD
                    \item Koszt wytworzenia po przliczeniu stawek za roboczo godziny: 2413456 USD
                    \item Koszt kontroli jakości po przliczeniu stawek za roboczo godziny: 3612820 USD
                \end{itemize}
            \item Rok 1986:
                \begin{itemize}
                    \item Wsparcie prac badawczo rozwojowych: 2718297 USD
                    \item Cena silnika i awioniki: 50976528 USD
                    \item Materiały: 57448933 USD
                    \item Próby w locie: 623756360 USD
                    \item Koszt inzynierów po przliczeniu stawek za roboczo godziny: 23053374 USD
                    \item Koszt narzędzi po przliczeniu stawek za roboczo godziny: 192492401 USD
                    \item Koszt wytworzenia po przliczeniu stawek za roboczo godziny: 20577415 USD
                    \item Koszt kontroli jakości po przliczeniu stawek za roboczo godziny: 20319358 USD
                \end{itemize}
        \end{itemize}

    \subsection{Podsumowanie analizy kosztów}

        \paragraph{}
            Podejrzewam, że autor \cite{galinski} nie bez powodu, wybrał datę 1970 i 1986.
            Chciał podkreślić to, iż o ile użycie oprogramowania CAD CAM CAE.
            Straca czas i zmniejsza koszty postania produktu.
        
        \paragraph{}
            W 1970 - używanie tych sytemów było bardzo drogie,
            więc większość używała desek reślarskich.
            Kadra inżynierska teżnie przyzwyczajona.
            W latach 70. nie było komputerów osobistych. 
            Systemy CAD działały na tzw. mainframe’ach (komputerach centralnych).
            Inżynierowie pracowali na terminalach graficznych, 
            które były niesamowicie drogie – jedno stanowisko mogło kosztować tyle, 
            co luksusowy dom.
        
        \paragraph{}
            W 1986 wtedy zaczęto używać tych narzędzi "dobrze",
            czyli w sposób zintegrowany.
            Modelowanie parametryczne (1987/1988):
            Firma PTC wypuściła program Pro/ENGINEER (dziś Creo).
            To był absolutny "game changer".
            Wprowadził on logikę,
            w której zmiana jednego wymiaru aktualizowała cały model

        \paragraph{}
            Zdecydowałem się sumowac kwotę z roku 1986 nastepnie przmnożyć ją przez wskażnik CPI.
            Zakładajac, iz roczna inflacja w okresie 2026-2030 2.5 \%. Szacowany kozt jednostkowy,
            przy zerowej marży przy założeniach sprzedarzy 1000 i produkcji 2 samolotów w miesiącu.
            Wynosi 3056193 USD.

        \paragraph{}
            Pytanie do promoamatora:
            \newline
            Czy methoda wynika z ramortu z lat 90 i skupupia sie głownie na produkcji smaolotow z durali.
            Czy do procesu produkcjnego drona też sie sprawdzi kwoty wydaja sie dośc wysokie.






